\documentclass[12pt]{article}
\usepackage{amsmath, amssymb, amsthm}
\usepackage{mathtools}
\usepackage{tensor}
\usepackage{geometry}
\usepackage{booktabs}
\usepackage{array}
\usepackage{xcolor}
\usepackage{hyperref}
\usepackage{fancyhdr}
\usepackage{slashed}
\geometry{margin=1.1in, headheight=15pt}

\definecolor{cadblue}{RGB}{30,90,200}
\definecolor{verifygreen}{RGB}{20,140,60}

\newcommand{\cdbexpr}[1]{\textcolor{cadblue}{$#1$}}
\newcommand{\verified}[1]{\textcolor{verifygreen}{\checkmark\; #1}}
\newcommand{\half}{\tfrac{1}{2}}
\newcommand{\dal}{\dot{\alpha}}
\newcommand{\dbe}{\dot{\beta}}
\newcommand{\dga}{\dot{\gamma}}
\newcommand{\dde}{\dot{\delta}}
\newcommand{\sigmabar}{\bar{\sigma}}
\newcommand{\omegap}{\omega_{\mathbf{p}}}

\pagestyle{fancy}
\fancyhf{}
\lhead{Srednicki QFT --- Chapter 37}
\rhead{Canonical Quantization of Spinor Fields I}
\cfoot{\thepage}

\title{\textbf{Srednicki QFT: Chapter 37}\\[6pt]
       \large Canonical Quantization of Spinor Fields I\\[4pt]
       \normalsize Cadabra2 expressions and numerical verification}
\author{Generated by \texttt{ch37\_export\_latex.py}}
\date{}

\begin{document}
\maketitle
\tableofcontents
\newpage

%% ============================================================
\section{Canonical Anticommutation Relations}
%% ============================================================

\subsection{Equal-time CARs}

The canonical anticommutation relations (CARs) for a left-handed Weyl field
$\psi_\alpha(x)$ are (Srednicki eq.~37.3):
\begin{equation}
  \boxed{
    \bigl\{\psi_\alpha(\mathbf{x},t),\;\psi^{\dagger\dot\beta}(\mathbf{y},t)\bigr\}
    = \sigma^0{}_{\alpha}{}^{\dot\beta}\,\delta^{(3)}(\mathbf{x}-\mathbf{y})
    = \delta_\alpha{}^{\dot\beta}\,\delta^{(3)}(\mathbf{x}-\mathbf{y})
  }
  \tag{37.3}
\end{equation}
where $\sigma^0 = I_2$ so the RHS is just $\delta_\alpha{}^{\dot\beta}$.  The remaining
equal-time anticommutators vanish:
\begin{align}
  \bigl\{\psi_\alpha(\mathbf{x},t),\,\psi_\beta(\mathbf{y},t)\bigr\} &= 0,
  \tag{37.4a}\\
  \bigl\{\psi^{\dagger\dot\alpha}(\mathbf{x},t),\,\psi^{\dagger\dot\beta}(\mathbf{y},t)\bigr\} &= 0.
  \tag{37.4b}
\end{align}

\textbf{Cadabra2 expression for the LHS:}
\[
  \psi_{\alpha} \psi^{\dagger \dal}\discretionary{}{}{}+\psi^{\dagger \dal} \psi_{\alpha} = \delta_{\alpha}\,^{\dal} \delta^{3\,}\left(x\discretionary{}{}{}-\,y\right)
\]

\subsection{Origin from Legendre transform}

The kinetic Lagrangian $\mathcal{L}_{\rm kin} = i\psi^\dagger \bar\sigma^\mu\partial_\mu\psi$
gives the conjugate momentum:
\begin{equation}
  \pi^\alpha = \frac{\partial\mathcal{L}}{\partial(\partial_0\psi_\alpha)}
  = i\psi^{\dagger\dot\beta}\,\bar\sigma^{0}{}_{\dot\beta}{}^\alpha
  = i\psi^{\dagger\alpha}
  \quad(\text{since }\bar\sigma^0 = I_2)
\end{equation}
Cadabra2: $\pi^\alpha = i \psidag_{\dal} \sigmabar^{0\, \dal \alpha} = i\psi^{\dagger\alpha}$.

The canonical equal-time anticommutation $\{\psi_\alpha(\mathbf{x}),\pi^\beta(\mathbf{y})\}
= i\delta_\alpha^\beta\,\delta^{(3)}(\mathbf{x}-\mathbf{y})$ then directly yields eq.~(37.3).

\medskip
\verified{$\sigma^0 = I_2$: max error $0.0e+00$.}

%% ============================================================
\section{Mode Expansion of the Weyl Field}
%% ============================================================

\subsection{Plane-wave expansion}

The left-handed Weyl field is expanded as (eq.~37.7):
\begin{equation}
  \psi_\alpha(x) = \int\!\frac{d^3p}{(2\pi)^3}\sum_s
  \Bigl[b_s(\mathbf{p})\,u^s_\alpha(\mathbf{p})\,e^{ip\cdot x}
       +d^\dagger_s(\mathbf{p})\,v^s_\alpha(\mathbf{p})\,e^{-ip\cdot x}\Bigr]
  \tag{37.7}
\end{equation}
with hermitian conjugate
\begin{equation}
  \psi^{\dagger\dot\alpha}(x) = \int\!\frac{d^3p}{(2\pi)^3}\sum_s
  \Bigl[b^\dagger_s(\mathbf{p})\,\tilde u^{\dot\alpha s}(\mathbf{p})\,e^{-ip\cdot x}
       +d_s(\mathbf{p})\,\tilde v^{\dot\alpha s}(\mathbf{p})\,e^{+ip\cdot x}\Bigr].
\end{equation}
Here $p\cdot x = p^\mu x_\mu$ with mostly-plus metric; on shell $p^0 = \omega_{\mathbf{p}} = \sqrt{\mathbf{p}^2+m^2}$.

\subsection{Operator CARs}

The creation/annihilation operators satisfy (eq.~37.8):
\begin{align}
  \{b_s(\mathbf{p}),\,b^\dagger_{s'}(\mathbf{p}')\}
  &= (2\pi)^3\,\delta^{(3)}(\mathbf{p}-\mathbf{p}')\,\delta_{ss'},
  \tag{37.8a}\\
  \{d_s(\mathbf{p}),\,d^\dagger_{s'}(\mathbf{p}')\}
  &= (2\pi)^3\,\delta^{(3)}(\mathbf{p}-\mathbf{p}')\,\delta_{ss'},
  \tag{37.8b}\\
  \{b_s,b_{s'}\} &= \{d_s,d_{s'}\} = \{b_s,d_{s'}\} = \cdots = 0.
  \tag{37.8c}
\end{align}

%% ============================================================
\section{Basis Spinors \texorpdfstring{$u^s(p)$}{u} and \texorpdfstring{$v^s(p)$}{v}}
%% ============================================================

\subsection{Weyl equations}

The basis spinors satisfy the Weyl equations (eq.~37.5):
\begin{align}
  p_{\alpha\dot\alpha}\,\tilde u^{\dot\alpha s}(p) &= +m\,u^s_\alpha(p),
  \tag{37.5a}\\
  p_{\alpha\dot\alpha}\,\tilde v^{\dot\alpha s}(p) &= -m\,v^s_\alpha(p),
  \tag{37.5b}
\end{align}
where $p_{\alpha\dot\alpha} = p_\mu\sigma^\mu_{\alpha\dot\alpha}$ is the momentum matrix in spinor space.

\subsection{Explicit construction at rest}

At $\mathbf{p}=0$, $p^\mu = (m,0,0,0)$, the momentum matrix is:
\[
  p_{\alpha\dot\alpha} = p_\mu\sigma^\mu_{\alpha\dot\alpha}
  = -m\,\sigma^0_{\alpha\dot\alpha} = -m\,I_2
  = \begin{pmatrix}-1 & 0 \\ 0 & -1\end{pmatrix}
\]

\textbf{Basis spinors at rest} (spin index $s = \pm$):
\begin{align*}
  u^+_\alpha &= \begin{pmatrix}1 \\ 0\end{pmatrix} & (\text{spin-up, particle})\\
  u^-_\alpha &= \begin{pmatrix}0 \\ 1\end{pmatrix} & (\text{spin-down, particle})\\
  v^+_\alpha &= \begin{pmatrix}0 \\ 1\end{pmatrix} & (\text{spin-up, antiparticle})\\
  v^-_\alpha &= \begin{pmatrix}-1 \\ 0\end{pmatrix} & (\text{spin-down, antiparticle})
\end{align*}

\subsection{Massless case}

For a massless particle with $p^\mu = (E,0,0,E)$:
\[
  p_{\alpha\dot\alpha}
  = -E\,\sigma^0 + E\,\sigma^3
  = \begin{pmatrix}0 & 0 \\ 0 & -2\end{pmatrix}
\]
Eigenvalues: $0, -2$.
The null eigenvector (helicity eigenstate $u^+$) lies in the kernel of
$p_{\alpha\dot\alpha}$, consistent with the massless Weyl equation
$p_{\alpha\dot\alpha}\tilde u^{\dot\alpha} = 0$.

\subsection{Completeness relations}

Summing over spins (eq.~37.10--37.11):
\begin{align}
  \sum_s u^s_\alpha(p)\,\tilde u^{\dot\beta s}(p) &= -p_{\alpha}{}^{\dot\beta} + m\,\delta_\alpha{}^{\dot\beta},
  \tag{37.10}\\
  \sum_s v^s_\alpha(p)\,\tilde v^{\dot\beta s}(p) &= -p_{\alpha}{}^{\dot\beta} - m\,\delta_\alpha{}^{\dot\beta}.
  \tag{37.11}
\end{align}

\textbf{Numerical check at rest:} $-p_{\alpha\dot\alpha} + m\,I_2 = 2m\,I_2$,
\[
  \sum_s u^s\otimes u^{s\dagger} = \begin{pmatrix}1 & 0 \\ 0 & 1\end{pmatrix}
  \approx \frac{1}{2m}\cdot \begin{pmatrix}2 & 0 \\ 0 & 2\end{pmatrix} \quad\checkmark
\]

\textbf{Vector bilinears at rest:}
\[
  \tilde u^{\dot\alpha +}\,\bar\sigma^\mu_{\dot\alpha\alpha}\,u^+_\alpha
  = [1.0, 0.0, 0.0, -1.0],\qquad
  \tilde u^{\dot\alpha -}\,\bar\sigma^\mu_{\dot\alpha\alpha}\,u^-_\alpha
  = [1.0, 0.0, 0.0, 1.0]
\]

%% ============================================================
\section{Deriving the CARs from the Mode Expansion}
%% ============================================================

Inserting eq.~(37.7) into $\{\psi_\alpha(x),\psi^{\dagger\dot\beta}(y)\}$ and
using the operator CARs (37.8a--b) (eq.~37.12):
\begin{align}
  \{\psi_\alpha(x),\psi^{\dagger\dot\beta}(y)\}
  &= \int\!\frac{d^3p}{(2\pi)^3}\sum_s u^s_\alpha(\mathbf{p})\,
     \tilde u^{\dot\beta s}(\mathbf{p})\,e^{ip(x-y)}
   + \int\!\frac{d^3p}{(2\pi)^3}\sum_s v^s_\alpha(\mathbf{p})\,
     \tilde v^{\dot\beta s}(\mathbf{p})\,e^{-ip(x-y)}
  \nonumber\\
  &= \int\!\frac{d^3p}{(2\pi)^3}\Bigl[
     \bigl(-p_\alpha{}^{\dot\beta}+m\delta_\alpha{}^{\dot\beta}\bigr)e^{ip(x-y)}
    +\bigl(-p_\alpha{}^{\dot\beta}-m\delta_\alpha{}^{\dot\beta}\bigr)e^{-ip(x-y)}
     \Bigr].
\end{align}
The momentum-dependent terms cancel between the two integrals, leaving:
\begin{equation}
  \{\psi_\alpha(x),\psi^{\dagger\dot\beta}(y)\}\Big|_{x^0=y^0}
  = \delta_\alpha{}^{\dot\beta}\,\delta^{(3)}(\mathbf{x}-\mathbf{y}).
  \qquad\checkmark
\end{equation}

%% ============================================================
\section{Normal-Ordered Hamiltonian}
%% ============================================================

\subsection{Classical Hamiltonian density}

From the Weyl Lagrangian via Legendre transform:
\begin{equation}
  H = \int\!d^3x\,\bigl[-i\psi^\dagger\bar\sigma^k\partial_k\psi
      + \tfrac{1}{2}m\psi\psi + \tfrac{1}{2}m\psi^\dagger\psi^\dagger\bigr].
\end{equation}

\subsection{Mode expansion form}

Inserting the mode expansion (eq.~37.17):
\begin{equation}
  H = \int\!\frac{d^3p}{(2\pi)^3}\,\omegap
  \Bigl[b^\dagger_s(\mathbf{p})b_s(\mathbf{p}) - d_s(\mathbf{p})d^\dagger_s(\mathbf{p})\Bigr].
\end{equation}

\subsection{Normal ordering}

Using $\{d_s(\mathbf{p}),d^\dagger_{s'}(\mathbf{p}')\} = (2\pi)^3\delta^{(3)}\delta_{ss'}$:
\begin{equation}
  -d_s\,d^\dagger_s = +d^\dagger_s\,d_s - \{d_s,d^\dagger_s\}
  = +d^\dagger_s\,d_s - (2\pi)^3\delta^{(3)}(0).
\end{equation}
Discarding the (divergent, negative) zero-point energy:
\begin{equation}
  \boxed{:H: = \int\!\frac{d^3p}{(2\pi)^3}\,\omegap
  \Bigl[b^\dagger_s(\mathbf{p})b_s(\mathbf{p}) + d^\dagger_s(\mathbf{p})d_s(\mathbf{p})\Bigr]}
  \tag{37.17}
\end{equation}
\textbf{Cadabra2:} $:H: = \displaystyle\int\!\frac{d^3p}{(2\pi)^3}\,\bigl(\omega_{p} b^{\dagger}\,_{s}\left(p\right) b_{s}\left(p\right)\discretionary{}{}{}+\omega_{p} d^{\dagger}\,_{s}\left(p\right) d_{s}\left(p\right)\bigr)$

\medskip
Key contrast with bosons:
\begin{center}
\renewcommand{\arraystretch}{1.4}
\begin{tabular}{@{}lll@{}}
  \toprule
  Field type & Zero-point energy & $:H:$ \\
  \midrule
  Boson (scalar) & $+\frac{\omega}{2}$ per mode & $\int\omega\,a^\dagger a$ \\
  Fermion (Weyl) & $-\frac{\omega}{2}$ per mode & $\int\omega\,(b^\dagger b + d^\dagger d)$ \\
  \bottomrule
\end{tabular}
\end{center}

%% ============================================================
\section{Spin-Statistics Theorem}
%% ============================================================

The spin-statistics theorem (Pauli 1940) states:
\begin{equation}
  \text{Spin-}\tfrac{1}{2} \implies \text{anticommuting quantization}.
\end{equation}
Proof sketch via microcausality: for a Lorentz-invariant theory,
\[
  \bigl[\mathcal{O}(x),\mathcal{O}'(y)\bigr] = 0
  \quad\text{for spacelike }(x-y)^2>0.
\]
If a spin-$\frac{1}{2}$ field were quantized with \emph{commutators}:
\[
  [\psi_\alpha(x),\psi^{\dagger\dot\beta}(y)] = \int\!\frac{d^3p}{(2\pi)^3}
  \bigl[(-p_{\alpha}{}^{\dot\beta}+m)e^{ip(x-y)}
       -(-p_{\alpha}{}^{\dot\beta}-m)e^{-ip(x-y)}\bigr]
  \neq 0\quad\text{(spacelike)}
\]
--- a violation.  With anticommutators the $p$-terms cancel and the result
vanishes outside the light cone, as required.

%% ============================================================
\section{Lorentz-Invariant Measure and Feynman Slash}
%% ============================================================

\subsection{The measure \texorpdfstring{$\widetilde{dp}$}{dp-tilde}}

The mode expansion uses the \textbf{Lorentz-invariant phase-space measure}
(Srednicki eq.~3.21):
\begin{equation}
  \widetilde{dp} \equiv \frac{d^3p}{(2\pi)^3\,2\omega_{\mathbf{p}}},
  \qquad \omega_{\mathbf{p}} = \sqrt{|\mathbf{p}|^2 + m^2}.
  \tag{3.21}
\end{equation}
It is manifestly Lorentz-invariant because it equals
$d^4p\,\delta(p^2+m^2)\,\theta(p^0)/(2\pi)^3$:
the 4-volume $d^4p$ is invariant, $\delta(p^2+m^2)$ restricts to the mass shell
(Lorentz scalar condition), and $\theta(p^0)$ selects the forward light-cone
(preserved by orthochronous boosts).

The mode expansion is then:
\begin{equation}
  \psi_\alpha(x) = \sum_s \int\!\widetilde{dp}\;
  \bigl[b_s(\mathbf{p})\,u^s_\alpha(\mathbf{p})\,e^{ip\cdot x}
       + d^\dagger_s(\mathbf{p})\,v^s_\alpha(\mathbf{p})\,e^{-ip\cdot x}\bigr].
\end{equation}

\subsection{Feynman slash notation}

The \textbf{Feynman slash} is defined by (Srednicki eq.~36.14):
\begin{equation}
  \slashed{p} \equiv p_\mu\gamma^\mu\,.
  \tag{36.14}
\end{equation}
(LaTeX: \verb|\slashed{p}| with \texttt{slashed} package, or \verb|\not{p}| built-in.)

\medskip
\textbf{Key property} (mostly-plus metric $g=\mathrm{diag}(-1,+1,+1,+1)$):
\begin{equation}
  \slashed{p}^2 = p_\mu p_\nu \gamma^\mu\gamma^\nu
  = \tfrac{1}{2}p_\mu p_\nu\{\gamma^\mu,\gamma^\nu\}
  = -p_\mu p^\mu = m^2 \quad(\text{on-shell})
  \tag{36.15}
\end{equation}

In the Weyl representation $\gamma^\mu = \bigl[\begin{smallmatrix}0&\sigma^\mu\\\bar\sigma^\mu&0\end{smallmatrix}\bigr]$,
the slash takes the block form:
\begin{equation}
  \slashed{p} = \begin{pmatrix}0 & p_{\alpha\dot\alpha} \\ \bar p^{\dot\alpha\alpha} & 0\end{pmatrix}
\end{equation}
The Dirac equation for plane-wave spinors (eq.~36.11 analog):
\begin{align}
  (\slashed{p} + m)\,u_s(p) &= 0, & (-\slashed{p}+m)\,v_s(p) &= 0, \tag{36.11}\\
  \bar u_s(p)\,(\slashed{p}+m) &= 0, & \bar v_s(p)\,(-\slashed{p}+m) &= 0.
\end{align}

\verified{$\slashed{p}^2 = m^2 I_4$ for on-shell $p$: max error $0.0e+00$.}

%% ============================================================
\section{Spin-Sum Completeness (4-component)}
%% ============================================================

The Dirac spin-sum completeness relations (Srednicki eqs.~38.28--38.29):
\begin{equation}
  \boxed{
    \sum_s u_s(p)\,\bar u_s(p) = -\slashed{p} + m\,,
    \qquad
    \sum_s v_s(p)\,\bar v_s(p) = -\slashed{p} - m\,.
  }
  \tag{38.28--29}
\end{equation}
where $\bar u = u^\dagger\gamma^0$ is the Dirac conjugate.

These are the 4-component version of the 2-component Weyl completeness relations
(eqs.~37.10--37.11).

\textbf{Physical use --- unpolarized cross sections:}  summing over unknown spins
replaces the outer product $u_s\bar u_{s'}$ with the spin-sum matrix, reducing
squared amplitudes to $\gamma$-matrix traces:
\begin{equation}
  \sum_{s,s'}|\bar u_{s'}(p')\,\Gamma\,u_s(p)|^2
  = \mathrm{Tr}\bigl[(-\slashed{p}'+m)\,\bar\Gamma\,(-\slashed{p}+m)\,\Gamma\bigr]
\end{equation}
where $\bar\Gamma = \gamma^0\Gamma^\dagger\gamma^0$.

\verified{$\sum_s u_s\bar u_s = -\slashed{p}+m$ at rest: max error $2.2e-16$.}
\verified{$\sum_s v_s\bar v_s = -\slashed{p}-m$ at rest: max error $2.2e-16$.}

%% ============================================================
\section{Summary}
%% ============================================================

\begin{center}
\renewcommand{\arraystretch}{1.5}
\begin{tabular}{@{}ll@{}}
  \toprule
  Object & Expression \\
  \midrule
  Equal-time CARs & $\{\psi_\alpha(\mathbf{x},t),\psi^{\dagger\dot\beta}(\mathbf{y},t)\}
    = \delta_\alpha{}^{\dot\beta}\delta^{(3)}(\mathbf{x}-\mathbf{y})$ \\
  Conjugate momentum & $\pi^\alpha = i\psi^{\dagger\alpha}$ \\
  Mode expansion & $\psi_\alpha = \int\frac{d^3p}{(2\pi)^3}\sum_s
    [b_s u^s_\alpha e^{ipx} + d^\dagger_s v^s_\alpha e^{-ipx}]$ \\
  Operator CARs & $\{b_s,b^\dagger_{s'}\} = \{d_s,d^\dagger_{s'}\}
    = (2\pi)^3\delta^{(3)}(\mathbf{p}-\mathbf{p}')\delta_{ss'}$ \\
  Weyl eq.\ (massive) & $p_{\alpha\dot\alpha}\tilde u^{\dot\alpha s} = +m u^s_\alpha$,\quad
    $p_{\alpha\dot\alpha}\tilde v^{\dot\alpha s} = -m v^s_\alpha$ \\
  Completeness & $\sum_s u^s_\alpha\tilde u^{\dot\beta s} = -p_\alpha{}^{\dot\beta} + m\delta_\alpha{}^{\dot\beta}$ \\
  Normal-ordered $H$ & $:H: = \int\frac{d^3p}{(2\pi)^3}\omegap[b^\dagger b + d^\dagger d]$ \\
  Spin-statistics & Spin-$\frac{1}{2}$ requires anticommuting fields (microcausality) \\
  Lorentz-invariant measure & $\widetilde{dp} = d^3p/[(2\pi)^3 2\omega_{\mathbf{p}}]$ \\
  Feynman slash & $\slashed{p}=p_\mu\gamma^\mu$;\quad $\slashed{p}^2=m^2$ on-shell \\
  Spin-sum (particle) & $\sum_s u_s\bar u_s = -\slashed{p}+m$ \\
  Spin-sum (antiparticle) & $\sum_s v_s\bar v_s = -\slashed{p}-m$ \\
  \bottomrule
\end{tabular}
\end{center}

\bigskip
\noindent\textbf{Next:} Chapter 38 develops the spinor algebra toolkit
($\sigma$-completeness, trace identities, Fierz rearrangement) needed for
explicit Feynman diagram calculations.

\end{document}